\documentclass[10pt,a4paper]{amsart}
\usepackage{amsmath,amssymb,amsthm,graphicx}
\usepackage{verbatim}
\newenvironment{reference}{\comment}{\endcomment}
%\newenvironment{reference}{}{}
\newenvironment{slogan}{\comment}{\endcomment}
\newenvironment{history}{\comment}{\endcomment}
% For commutative diagrams we use Xy-pic
\usepackage[all]{xy}

% We use 2cell for 2-commutative diagrams.
\xyoption{2cell}
% Theorem environments.
%
\theoremstyle{plain}
\newtheorem{theorem}[subsection]{Theorem}
\newtheorem{proposition}[subsection]{Proposition}
\newtheorem{lemma}[subsection]{Lemma}

\theoremstyle{definition}
\newtheorem{definition}[subsection]{Definition}
\newtheorem{example}[subsection]{Example}
\newtheorem{exercise}[subsection]{Exercise}
\newtheorem{situation}[subsection]{Situation}

\theoremstyle{remark}
\newtheorem{remark}[subsection]{Remark}
\newtheorem{remarks}[subsection]{Remarks}

\numberwithin{equation}{subsection}

% Macros
%
\def\lim{\mathop{\mathrm{lim}}\nolimits}
\def\colim{\mathop{\mathrm{colim}}\nolimits}
\def\Spec{\mathop{\mathrm{Spec}}}
\def\Hom{\mathop{\mathrm{Hom}}\nolimits}
\def\Ext{\mathop{\mathrm{Ext}}\nolimits}
\def\SheafHom{\mathop{\mathcal{H}\!\mathit{om}}\nolimits}
\def\SheafExt{\mathop{\mathcal{E}\!\mathit{xt}}\nolimits}
\def\Sch{\mathit{Sch}}
\def\Mor{\mathop{\mathrm{Mor}}\nolimits}
\def\Ob{\mathop{\mathrm{Ob}}\nolimits}
\def\Sh{\mathop{\mathit{Sh}}\nolimits}
\def\NL{\mathop{N\!L}\nolimits}
\def\CH{\mathop{\mathrm{CH}}\nolimits}
\def\proetale{{pro\text{-}\acute{e}tale}}
\def\etale{{\acute{e}tale}}
\def\QCoh{\mathit{QCoh}}
\def\Ker{\mathop{\mathrm{Ker}}}
\def\Im{\mathop{\mathrm{Im}}}
\def\Coker{\mathop{\mathrm{Coker}}}
\def\Coim{\mathop{\mathrm{Coim}}}

% Boxtimes
%
\DeclareMathSymbol{\boxtimes}{\mathbin}{AMSa}{"02}

%
% Macros for moduli stacks/spaces
%
\def\QCohstack{\mathcal{QC}\!\mathit{oh}}
\def\Cohstack{\mathcal{C}\!\mathit{oh}}
\def\Spacesstack{\mathcal{S}\!\mathit{paces}}
\def\Quotfunctor{\mathrm{Quot}}
\def\Hilbfunctor{\mathrm{Hilb}}
\def\Curvesstack{\mathcal{C}\!\mathit{urves}}
\def\Polarizedstack{\mathcal{P}\!\mathit{olarized}}
\def\Complexesstack{\mathcal{C}\!\mathit{omplexes}}
% \Pic is the operator that assigns to X its picard group, usage \Pic(X)
% \Picardstack_{X/B} denotes the Picard stack of X over B
% \Picardfunctor_{X/B} denotes the Picard functor of X over B
\def\Pic{\mathop{\mathrm{Pic}}\nolimits}
\def\Picardstack{\mathcal{P}\!\mathit{ic}}
\def\Picardfunctor{\mathrm{Pic}}
\def\Deformationcategory{\mathcal{D}\!\mathit{ef}}
\usepackage{cleveref}
\crefname{originalreference}{original reference}{original references}
\Crefname{originalreference}{Original reference}{Original references}
\makeatletter
\newlabel{divisors-lemma-blowing-up-affine}{{divisors-lemma-blowing-up-affine}{source}{}{sourcepointer-1}{}}
\newlabel{divisors-lemma-blowing-up-affine@cref}{{[originalreference][0][]divisors-lemma-blowing-up-affine}{[1][0][]source}{}{}{}}
\makeatother
\crefname{originallemma}{original lemma}{original lemmas}
\Crefname{originallemma}{Original lemma}{Original lemmas}
\makeatletter
\newlabel{lemma-factor-hom}{{lemma-factor-hom}{source}{}{sourcepointer-2}{}}
\newlabel{lemma-factor-hom@cref}{{[originallemma][0][]lemma-factor-hom}{[1][0][]source}{}{}{}}
\makeatother
\crefname{originalreference}{original reference}{original references}
\Crefname{originalreference}{Original reference}{Original references}
\makeatletter
\newlabel{more-algebra-lemma-ext-1-annihilated-definite}{{more-algebra-lemma-ext-1-annihilated-definite}{source}{}{sourcepointer-3}{}}
\newlabel{more-algebra-lemma-ext-1-annihilated-definite@cref}{{[originalreference][0][]more-algebra-lemma-ext-1-annihilated-definite}{[1][0][]source}{}{}{}}
\makeatother
\crefname{originalreference}{original reference}{original references}
\Crefname{originalreference}{Original reference}{Original references}
\makeatletter
\newlabel{more-algebra-lemma-map-out-of-almost-free}{{more-algebra-lemma-map-out-of-almost-free}{source}{}{sourcepointer-4}{}}
\newlabel{more-algebra-lemma-map-out-of-almost-free@cref}{{[originalreference][0][]more-algebra-lemma-map-out-of-almost-free}{[1][0][]source}{}{}{}}
\makeatother
\title{The comparison map on Ext of a two-term complex over the quotient vanishes after a shift}
\author{The Stacks Project Authors}
\date{stacks-project:local-cohomology.tex:local-snapshot-20261012; GNU FDL 1.2 or later}
\begin{document}
\maketitle
\noindent\textit{Editorial note.} All delivered mathematical text is the author's own, copied from the pinned Stacks Project chapter source sources/stacks-local-cohomology/pinned-original-source.tex; no statement or proof step was rewritten, reordered or removed.\par\medskip
\noindent\textit{Editorial note.} Preamble provenance: the Stacks Project keeps one single shared preamble for every chapter; this card takes it directly from the staged sources/stacks-local-cohomology/preamble.tex as three verbatim spans (lines 8-12, 14-18 and 154-217), so no derived file was written and nothing was added or redefined. The wrapper class and the base packages are supplied by the assembler.\par\medskip
\noindent\textit{Editorial note.} Bibliography provenance: neither the counted statement nor its proof contains a citation (this lane rescans the whole span before writing the card), so the card carries no bibliography and no entry was reconstructed.\par\medskip
\noindent\textit{Editorial note.} Selection provenance: this card was opened under the file-based control contract (CONTROL.md), with the mechanical pre-card gate run on its statement and proof interval before the card was written.\par\medskip
\noindent\textit{Allowed standard prerequisites (editorial).}\begin{itemize}
\item The representation of a bounded complex over the quotient ring by a two-term complex whose first term is projective for the ideal power. Invoked inside the counted proof at source lines 5469-5476 and attributed to a lemma of another chapter (a cross-chapter pointer). Provides the explicit complex on which the comparison map is computed.
\item The description of maps out of an almost free module and the factorisation lemma of this chapter. Both are invoked inside the counted proof (a cross-chapter pointer and the author's own lemma lemma-factor-hom). Convert the required vanishing of the comparison map into a factorisation statement, which is then verified diagrammatically.
\end{itemize}\medskip
\ifdefined\hypertarget\hypertarget{sourcepointer-1}{}\fi
\noindent\textit{Source reference.} Original reference divisors-lemma-blowing-up-affine: sources/stacks-local-cohomology/pinned-original-source.tex, source lines 4655--4655. Navigation pointer only: the counted text refers to the Stacks item divisors-lemma-blowing-up-affine, which belongs to another chapter of the project and is not part of this card. The pointer carries the target's own label as its number because this chapter's staged snapshot does not contain the other chapter's tag table; its content is not reproduced here.\par\medskip
\ifdefined\hypertarget\hypertarget{sourcepointer-2}{}\fi
\noindent\textit{Source reference.} Original lemma lemma-factor-hom: sources/stacks-local-cohomology/pinned-original-source.tex, source lines 5486--5486. Navigation pointer only: the counted text refers to the Stacks item lemma-factor-hom, which belongs to another chapter of the project and is not part of this card. The pointer carries the target's own label as its number because this chapter's staged snapshot does not contain the other chapter's tag table; its content is not reproduced here.\par\medskip
\ifdefined\hypertarget\hypertarget{sourcepointer-3}{}\fi
\noindent\textit{Source reference.} Original reference more-algebra-lemma-ext-1-annihilated-definite: sources/stacks-local-cohomology/pinned-original-source.tex, source lines 5471--5471. Navigation pointer only: the counted text refers to the Stacks item more-algebra-lemma-ext-1-annihilated-definite, which belongs to another chapter of the project and is not part of this card. The pointer carries the target's own label as its number because this chapter's staged snapshot does not contain the other chapter's tag table; its content is not reproduced here.\par\medskip
\ifdefined\hypertarget\hypertarget{sourcepointer-4}{}\fi
\noindent\textit{Source reference.} Original reference more-algebra-lemma-map-out-of-almost-free: sources/stacks-local-cohomology/pinned-original-source.tex, source lines 5475--5475. Navigation pointer only: the counted text refers to the Stacks item more-algebra-lemma-map-out-of-almost-free, which belongs to another chapter of the project and is not part of this card. The pointer carries the target's own label as its number because this chapter's staged snapshot does not contain the other chapter's tag table; its content is not reproduced here.\par\medskip
\begin{lemma}
\label{lemma-bound-two-term-complex}
With $d = d(S)$ and $q_0 = q(S)$ as above. Then
\begin{enumerate}
\item for integers $n \geq c \geq 0$ with
$n \geq \max(q_0 + (1 + d)c, (2 + d)c)$,
\item for $K$ of $D(A/I^n)$ with $H^i(K) = 0$ for $i \not = -1, 0$
and $H^i(K)$ finite for $i = -1, 0$ such that $\Ext^1_{A/I^c}(K, N)$
is annihilated by $I^c$ for all finite $A/I^n$-modules $N$
\end{enumerate}
the map
$$
\Ext^1_{A/I^n}(K, I^n/I^{2n})
\longrightarrow
\Ext^1_{A/I^n}(K, I^{n - (1 + d)c}/I^{2n - 2(1 + d)c})
$$
is zero.
\end{lemma}

\begin{proof}
The case $d > 0$. Let $K^{-1} \to K^0$ be a complex representing $K$ as in
More on Algebra, Lemma \ref{more-algebra-lemma-ext-1-annihilated-definite}
part (5) with respect to the ideal $I^c/I^n$ in the ring $A/I^n$.
In particular $K^{-1}$ is $I^c/I^n$-projective as multiplication
by elements of $I^c/I^n$ even factor through $K^0$. By
More on Algebra, Lemma \ref{more-algebra-lemma-map-out-of-almost-free} part (1)
we have
$$
\Ext^1_{A/I^n}(K, I^n/I^{2n}) =
\Coker(\Hom_{A/I^n}(K^0, I^n/I^{2n}) \to \Hom_{A/I^n}(K^{-1}, I^n/I^{2n}))
$$
and similarly for other Ext groups. Hence any class $\xi$ in
$\Ext^1_{A/I^n}(K, I^n/I^{2n})$
comes from an element $\varphi \in \Hom_{A/I^n}(K^{-1}, I^n/I^{2n})$.
Denote $\varphi'$ the image of $\varphi$ in
$\Hom_{A/I^n}(K^{-1}, I^{n - (1 + d)c}/I^{2n - (1 + d)c})$.
By Lemma \ref{lemma-factor-hom}
we can write $\varphi' = \sum a_i \psi_i$ with $a_i \in I^c$ and
$\psi_i \in \Hom_{A/I^n}(M, I^{n - (2 + d)c}/I^{2n - (2 + d)c})$.
Choose $h_i : K^0 \to K^{-1}$ such that
$a_i \text{id}_{K^{-1}} = h_i \circ d_K^{-1}$. Set
$\psi = \sum \psi_i \circ h_i : K^0 \to I^{n - (2 + d)c}/I^{2n - (2 + d)c}$.
Then $\varphi' = \psi \circ \text{d}_K^{-1}$ and we conclude that
$\xi$ already maps to zero in
$\Ext^1_{A/I^n}(K, I^{n - (1 + d)c}/I^{2n - (1 + d)c})$
and a fortiori in
$\Ext^1_{A/I^n}(K, I^{n - (1 + d)c}/I^{2n - 2(1 + d)c})$.

\medskip\noindent
The case $d = 0$\footnote{The argument given for $d > 0$ works but gives
a slightly weaker result.}. Let $\xi$ and $\varphi$ be as above.
We consider the diagram
$$
\xymatrix{
K^0 \\
K^{-1} \ar[u] \ar[r]^\varphi & I^n/I^{2n} \ar[r] & I^{n - c}/I^{2n - c}
}
$$
Pulling back to $X$ and using the map
$p^*(I^n/I^{2n}) \to \mathcal{O}_{Y_n}(n)$
we find a solid diagram
$$
\xymatrix{
p^*K^0 \ar@{..>}[rrd] \\
p^*K^{-1} \ar[u] \ar[r] &
\mathcal{O}_{Y_n}(n) \ar[r] &
\mathcal{O}_{Y_n}(n - c)
}
$$
We can cover $X$ by affine opens $U = \Spec(B)$ such that there
exists an $a \in I$ with the following property: $IB = aB$
and $a$ is a nonzerodivisor on $B$. Namely, we can cover
$X$ by spectra of affine blowup algebras, see
Divisors, Lemma \ref{divisors-lemma-blowing-up-affine}.
The restriction of $\mathcal{O}_{Y_n}(n) \to \mathcal{O}_{Y_n}(n - c)$
to $U$ is isomorphic to the map of quasi-coherent $\mathcal{O}_U$-modules
corresponding to the $B$-module map $a^c : B/a^nB \to B/a^nB$.
Since $a^c : K^{-1} \to K^{-1}$ factors through $K^0$ we see that
the dotted arrow exists over $U$. In other words,
locally on $X$ we can find the dotted arrow! Now the sheaf of dotted
arrows fitting into the diagram is principal homogeneous under
$$
\mathcal{F} = \SheafHom_{\mathcal{O}_X}(
\Coker(p^*K^{-1} \to p^*K^0), \mathcal{O}_{Y_n}(n - c))
$$
which is a coherent $\mathcal{O}_X$-module.
Hence the obstruction for finding the dotted arrow is an element
of $H^1(X, \mathcal{F})$. This cohomology group is zero as
$1 > d = 0$, see discussion following the definition of $d = d(S)$.
This proves that we can find a dotted arrow
$\psi : p^*K^0 \to \mathcal{O}_{Y_n}(n - c)$
fitting into the diagram. Since $n - c \geq q_0$
we find that $\psi$ induces a map $K^0 \to I^{n - c}/I^{2n - c}$.
Chasing the diagram we conclude that $\varphi' = \psi \circ \text{d}_K^{-1}$
and the proof is finished as before.
\end{proof}

\end{document}
